\partnote{Symmetry and effective descent}{A set of solutions forgets the witnesses relating equivalent solutions. The next layer retains those witnesses, their composition, and their automorphisms. Gluing then has to respect the overlap identifications themselves.}
\section{Nonabelian coherence and relative trivialization}\label{sec:gauge}
\subsection{Coherent transition labels}
Fix groups \(G_v\) and restriction homomorphisms \(f^*:G_w\to G_v\), with the same identity and composition laws as before. The groups need not commute and the homomorphisms need not be bijective.

\begin{definition}[Cocycle]
A cocycle is a family \(c_f\in G_v\) for every \(f:v\to w\), satisfying
\[
 c_{1_v}=1,\qquad c_{g\circ f}=c_f f^*(c_g).
\]
It determines a twisted view system with maps
\[
 T_c(f)(x)=c_f f^*(x).
\]
A section of the twisted system satisfies \(x_v=c_f f^*(x_w)\).
\end{definition}
The cocycle law says that the transition rules compose coherently. It does not assert that the twisted system has a section. The coherent transition system and a global trivialization of that system are different objects.

For a family \(a=(a_v)\in\prod_vG_v\), define
\begin{equation}\label{eq:gauge}
 (a\cdot c)_f=a_vc_f[f^*(a_w)]^{-1}.
\end{equation}

\begin{proposition}[Gauge action and its groupoid]\label{KS30}
Equation~\eqref{eq:gauge} defines an action on cocycles. A gauge carrying \(c\) to \(d\) induces an equivalence of twisted section sets by \(x_v\mapsto a_vx_v\). The cocycles and actual gauges form a groupoid: an arrow \(a:c\to d\) followed by \(b:d\to e\) has label \(ba\).
\end{proposition}
\begin{proof}
For composable \(f:v\to w\) and \(g:w\to z\), the middle factors cancel in
\[
 (a\cdot c)_f f^*((a\cdot c)_g)
 =a_vc_f f^*(c_g)[f^*g^*(a_z)]^{-1}.
\]
This is the transformed label at \(g\circ f\). The identity law is similar. Direct substitution gives \(a\cdot(b\cdot c)=(ab)\cdot c\), with identity gauge acting trivially. If \(x\) is twisted coherent, then
\[
 (a\cdot c)_f f^*(a_wx_w)=a_vc_f f^*(x_w)=a_vx_v.
\]
The inverse gauge gives the inverse section map. Identity arrows, composition, and inverses in the groupoid now follow from the group action laws.
\end{proof}
\source{KS30}
The action groupoid \(Z^1(\V;G)//\prod_vG_v\) retains comparison witnesses. Its orbit set \(Z^1(\V;G)/\prod_vG_v\) retains only isomorphism classes. There is no reason for the orbit set to be a group in this nonabelian setting.

\subsection{Automorphisms and ambiguity are not the same question}
The stabilizer of a cocycle is the group of families with
\[
 a_v=c_f f^*(a_w)c_f^{-1}.
\]
If one gauge \(a:c\to d\) is known, all gauges from \(c\) to \(d\) are uniquely of the form \(ah\), with \(h\in\Stab(c)\). Indeed \(a^{-1}b\) stabilizes \(c\) for every other gauge \(b\). This is the ordinary Hom torsor principle in a groupoid.

That symmetry group must not be confused with the set of global trivializations of one fixed cocycle. The latter will be parametrized by ordinary coherent group elements, after one trivialization has been selected.

\subsection{The relative obstruction without commutativity}
Prescribe a coherent partial section \(s\) of \(T_c\) on \(U\). The based gauge group is
\[
 \mathcal G_U=\{a\in\prod_vG_v:a_u=1\text{ for }u\in U\}.
\]
Choose a raw family \(q\) equal to \(s\) on \(U\). Identity elements can be used elsewhere. Normalize the labels by
\[
 c^q_f=q_v^{-1}c_f f^*(q_w).
\]
Its labels on arrows internal to \(U\) are one. Define the pointed orbit obstruction
\[
 \mathfrak o_U(c,s)=[c^q]\in Z^1(\V;G)/\mathcal G_U,
\]
whose neutral point is the class of the trivial cocycle.

\begin{theorem}[Nonabelian relative extension]\label{KS31}
The obstruction is independent of the raw lift. It is neutral exactly when \(s\) extends to a global section of \(T_c\).
\end{theorem}
\begin{proof}
If \(t\) is another raw lift, the family \(t^{-1}q\) is based and carries \(c^q\) to \(c^t\). Thus the orbit does not depend on the lift.

If a based gauge \(b\) carries \(c^q\) to the trivial cocycle, set \(x_v=q_vb_v^{-1}\). The gauge equation becomes \(x_v=c_f f^*(x_w)\), and \(b|_U=1\) gives \(x|_U=s\). Conversely, an extension \(x\) gives the based gauge \(b=x^{-1}q\), which trivializes \(c^q\). All products retain the displayed order. No commutativity has been used.
\end{proof}
\source{KS31}
The orbit is defined from actual cocycles and an actual group action. It is not a Boolean test defined to be neutral precisely when an extension exists. Deciding orbit membership in arbitrary groups remains a separate problem.

Let
\[
 H_U=\{h:f^*(h_w)=h_v\text{ for every }f,\ h|_U=1\}.
\]

\begin{theorem}[The nonabelian extension torsor]\label{KS32}
Given one extension \(x_0\), the map \(h\mapsto x_0h\) is an equivalence from \(H_U\) to the extension set. The extension is unique if and only if the obstruction is neutral and \(H_U\) is trivial.
\end{theorem}
\begin{proof}
If \(h\) is coherent, then
\[
 c_f f^*((x_0)_wh_w)=(x_0)_vh_v,
\]
so right multiplication gives an extension. If \(x\) is another extension, then
\[
 (x_0)_v^{-1}x_v
 =[c_f f^*((x_0)_w)]^{-1}[c_f f^*(x_w)]
 =f^*((x_0)_w^{-1}x_w).
\]
Thus \(x_0^{-1}x\in H_U\) and is the inverse parameter. The uniqueness statement follows as in the additive case.
\end{proof}
\source{KS32}
The additive theory is recovered by applying a multiplicative type tag to an additive view system and choosing the trivial cocycle. The normalized label then becomes \(f^*q_w-q_v\). The source proves agreement of the actual extension sets and vanishing criteria. It does not identify the whole nonabelian orbit set with every element of the larger additive arrow residual quotient.

\subsection{A small noncommutative model}
\begin{example}[A \(C_2\) loop with \(S_3\) coefficients]\label{KS33}
Take one view with a loop \(\sigma\) satisfying \(\sigma^2=1\), and constant coefficients in a group \(H\). A cocycle is determined by an involution \(p\in H\). Gauge equivalence is conjugacy, and its automorphisms form the centralizer of \(p\).

For \(H=S_3\), the involutions are the identity and the three transpositions. There are exactly two orbit classes. The identity cocycle has automorphism group \(S_3\), which is noncommutative. A nonidentity involution gives a coherent cocycle whose regular twisted section equation \(x=px\) has no solution.
\end{example}
\begin{proof}
The only nontrivial cocycle relation is \(p^2=1\). At a single view the gauge action is \(p\mapsto apa^{-1}\). All three transpositions are conjugate and none is conjugate to the identity. The stabilizer is therefore the centralizer. Finally \(x=px\) implies \(p=1\) after multiplication by \(x^{-1}\).
\end{proof}
\source{KS33}
The example makes an important boundary visible. A cocycle can exist and glue as an object even when it admits no global trivializing section. Stackhood will concern the former gluing problem, not a claim that every cocycle becomes trivial.

\plate{06_groupoid}{The orbit set sees two classes in the \(C_2\)-to-\(S_3\) example. The groupoid also retains the arrows between conjugate transpositions and all automorphisms. Collapsing the arrows loses information that matters in descent.}{fig:groupoid}

\chap{Effective gluing of objects and their identifications}\label{sec:descent}
\subsection{Why pairwise agreement is not the right datum}
For objects with symmetry, local descriptions need not be literally equal on overlaps. They may instead come with specified isomorphisms. Those isomorphisms must satisfy their own composition law. It is not enough to know that some isomorphism exists for each pair.

Let \(U_i\) be full patches in \(\V\). Assume the composition cover condition
\begin{equation}\label{eq:compositioncover}
 \forall(v\xrightarrow f w\xrightarrow g z),\quad
 \exists i:\ v,w,z\in U_i.
\end{equation}
By inserting identity arrows, this implies coverage of objects and arrows.

A descent datum consists of local cocycles \(c_i\) and overlap gauges \(a_{ij}(v)\) with
\begin{align}
 c_j(f)&=a_{ij}(v)c_i(f)[f^*a_{ij}(w)]^{-1},\label{eq:localgauge}\\
 a_{ik}(v)&=a_{jk}(v)a_{ij}(v).\label{eq:triple}
\end{align}
The triple equation implies \(a_{ii}=1\) and \(a_{ji}=a_{ij}^{-1}\). A morphism of descent data is a family of local gauges \(r_i\) satisfying
\[
 r_j(v)a_{ij}(v)=a'_{ij}(v)r_i(v).
\]
These objects and arrows form a groupoid.

\begin{theorem}[Effective cocycle descent]\label{KS34}
Under \eqref{eq:compositioncover}, restriction is an equivalence between the global cocycle gauge groupoid and the groupoid of descent data. The overlap gauges supplied in the datum are realized, not discarded.
\end{theorem}
\begin{proof}
Choose a reference patch \(\kappa(v)\) through each view and set
\[
 b_i(v)=a_{\kappa(v),i}(v).
\]
Equation~\eqref{eq:triple} gives \(b_j=a_{ij}b_i\). Normalize the local labels by
\[
 \widetilde c_i(f)=b_i(v)^{-1}c_i(f)f^*(b_i(w)).
\]
Substituting \eqref{eq:localgauge} and \(b_j=a_{ij}b_i\) shows that these normalized labels agree literally on overlaps. They can therefore be assembled into labels \(c_f\). For every composable pair, a patch containing all three endpoints verifies the cocycle law for the assembled labels. Thus \(c\) is a global cocycle.

Undoing the normalization gives
\[
 b_i(v)c_f[f^*(b_i(w))]^{-1}=c_i(f),\qquad b_j=a_{ij}b_i.
\]
So the local identifications implement the original overlap gauges. This constructs an object whose restriction is isomorphic to the specified descent datum.

For morphisms between restrictions of two global cocycles, the overlap gauges are identities. The morphism condition reduces to equality of the local gauges on overlaps. They assemble uniquely into a global gauge, and arrow coverage verifies its compatibility with the cocycles. Restriction and assembly are inverse on Hom sets. The restriction functor is therefore fully faithful and essentially surjective, hence an equivalence of groupoids.
\end{proof}
\source{KS34}

\subsection{The correct uniqueness statement}
Suppose \((c,b_i)\) and \((c',b'_i)\) realize the same descent datum. The only possible compatible comparison is
\[
 g_v=b'_i(v)^{-1}b_i(v),
\]
using any patch through \(v\). The overlap identities make it independent of the patch, and \(b'_ig=b_i\). It is therefore the unique global comparison respecting the fixed local identifications.

This is not a statement that the underlying cocycle has no automorphisms. Those automorphisms can reappear when the local identifications are forgotten. Nor does the construction choose one representative independently of every reference patch. It proves that the choices are related by the appropriate unique comparison.

\subsection{A failed triple}
Let \(t\ne1\) be an involution and ask for frames \(b_0,b_1,b_2\) satisfying
\[
 b_1=tb_0,\qquad b_2=tb_1,\qquad b_2=tb_0.
\]
The first two give \(b_2=b_0\), while the third forces \(t=1\). Each displayed transition is individually invertible, but they cannot be the overlap identifications of a global frame. The defect lies in the triple relation.

\plate{07_cover}{Object coverage, arrow coverage, and composition coverage are distinct. The proper two-view patches on a three-view chain see every arrow but no patch sees the relation \(0\to1\to2\). The four proper three-view patches on a four-view chain do cover every composable pair.}{fig:cover}

\chap{A genuine stack on the composition site}\label{sec:stack}
The preceding theorem can be placed in a standard site, rather than called stack-like. We use the usual effective-descent meaning of a stack~\cite{stacks-stack,mathlib-stack}. The base of this site is the inclusion poset of full patches
\[
 \mathcal P(\V)=(\operatorname{Set}(\Ob\V),\subseteq).
\]
This patch category is different from the original category of views.

\subsection{The topology is defined without assuming effectivity}
For a composable pair \(q=(v\xrightarrow f w\xrightarrow g z)\), let its support be \(|q|=\{v,w,z\}\). A sieve \(\mathcal R\) on \(W\) is covering when
\begin{equation}\label{eq:site}
 |q|\subseteq W\quad\Longrightarrow\quad
 (|q|\hookrightarrow W)\in\mathcal R
 \quad\text{for every composable }q.
\end{equation}
This condition depends only on the view category, not on the coefficient groups or on a promised global realization.

\begin{proposition}[The composition topology]\label{KS35}
Condition~\eqref{eq:site} defines a Grothendieck topology \(J_{\mathrm{comp}}\). The sieve generated by a family of inclusions \(U_i\subseteq W\) is covering exactly when the family covers every composition support in \(W\).
\end{proposition}
\begin{proof}
The maximal sieve contains all support inclusions. If \(Z\subseteq W\), every support in \(Z\) is also a support in \(W\), so its inclusion belongs to the pullback of a covering sieve. For transitivity, let \(\mathcal R\) cover \(W\), and let another sieve \(\mathcal T\) cover locally over every member of \(\mathcal R\). The support inclusion \(|q|\hookrightarrow W\) belongs to \(\mathcal R\). Applying the local covering condition on \(|q|\) to the same composable pair puts the identity inclusion of \(|q|\) into the pulled back \(\mathcal T\). Thus \(|q|\hookrightarrow W\) belongs to \(\mathcal T\).

In the poset of subsets, an inclusion of a support factors through a member of the family precisely when that member contains the support. Membership in the generated sieve is exactly this factorization. This proves the final characterization.
\end{proof}
\source{KS35}
Intersections of a covering family with a subpatch remain covering. Replacing each member by a covering family again gives a cover. Identities imply coverage of single views and arrows. The empty patch admits the empty cover.

\subsection{The actual patch pseudofunctor}
On a patch \(W\), use
\[
 \mathfrak K_G(W)=Z^1(\V|_W;G|_W)//\prod_{v\in W}G_v.
\]
Restriction keeps the existing cocycle values and gauge values on the smaller patch. It first gives a strict category valued functor on the opposite patch poset. Its standard pseudofunctor form has identity comparison isomorphisms, with the required coherence laws.

\begin{theorem}[Stackhood on the composition site]\label{KS36}
The patch pseudofunctor \(\mathfrak K_G\) is a stack for \(J_{\mathrm{comp}}\). It satisfies the standard condition of descent for morphisms and effective descent for objects on every covering sieve.
\end{theorem}
\begin{proof}
For a covering family, Theorem~\ref{KS34} gives a concrete equivalence of groupoids. The standard descent interface allows comparison morphisms on arbitrary common subpatches, not just pointwise labels. Evaluate these comparison morphisms on singleton patches. Their naturality under restriction identifies those values with their values on every larger common subpatch. The standard cocycle law becomes \eqref{eq:triple}. Conversely, the concrete effectuation gives the required standard local identifications, and the same naturality proves every common subpatch equation.

The morphism construction retains all Hom sets and their inverse gluing maps. Thus the standard restriction functor is fully faithful and essentially surjective. Finally, index an arbitrary covering sieve by its own arrows. Every composition support is one of those arrows by definition, so that family satisfies \eqref{eq:compositioncover}. The family equivalence proves effective descent for every covering sieve, which is the standard stack condition.
\end{proof}
\source{KS36}
The principal source declaration is \code{GroupViewSystem.patch_isStack}. Its target is mathlib's \code{Pseudofunctor.IsStack}, not a project predicate that repeats the conclusion as an assumption.

\subsection{The cover is nontrivial, but specific}
On a four-view chain, take the four patches obtained by deleting one vertex. All patches are proper, yet every composable pair has at most three endpoints and lies in one of them. The theorem genuinely glues pieces smaller than the whole object.

On the three-view chain, the analogous proper patches cover each object and each arrow but miss the support of \(0\to1\to2\). They do not form a cover in this topology. Consequently the theorem is neither vacuous nor an unrestricted gluing principle.

\begin{scopebox}
\textbf{Exact scope.} This is an ordinary stack of strict group valued cocycle realizations on a specified combinatorial site. It is not an infinity-stack, a universal heterogeneous atlas, or an arithmetic existence theorem. The stack theorem glues coherent cocycles with their overlap isomorphisms. It does not assert that every resulting cocycle has a global trivialization.
\end{scopebox}

\chap{Equivalent presentations and identified boundaries}\label{sec:relative}
\subsection{Transporting the same question}
Let \(H:\V\to\W\) be part of a category equivalence, with inverse \(L\) and counit \(\varepsilon:HL\Rightarrow\id_\W\). Pull a view system \(F\) on \(\W\) back along \(H\). Sampling a global section gives \((H^*x)_v=x_{Hv}\).

\begin{proposition}[Equivalence of section spaces]\label{KS37}
Pullback induces \(\Gamma(F)\simeq\Gamma(H^*F)\). If \(U\subseteq\Ob\W\) is closed under isomorphism and \(U'=H^{-1}U\), it also induces
\[
 \Ext_U^F(s)\simeq\Ext_{U'}^{H^*F}(H^*s).
\]
Existence and uniqueness of the prescribed extension are preserved and reflected.
\end{proposition}
\begin{proof}
For a section \(y\) of the pulled back system, use \(t_w=\varepsilon_w^{-1}:w\to HLw\) and set
\[
 (Ry)_w=F(t_w)(y_{Lw}).
\]
Naturality and coherence prove that this is a section. The unit and counit identities show that restriction and recovery are inverse.

For a prescribed domain, \(w\in U\) implies \(HLw\in U\) because these objects are isomorphic. Thus the partial equation for \(y\) applies at \(Lw\). Coherence of \(s\) along \(t_w\) gives \((Ry)_w=s_w\). Hence the same inverse maps restrict to the extension fibers. Any equivalence of sets preserves and reflects nonemptiness and the singleton property.
\end{proof}
\source{KS37}
The coefficients really are pulled back. This does not compare unrelated coefficient systems simply because their diagrams look alike. Isomorphism closure is a sufficient general hypothesis for the relative statement. The theorem does not assert that every unsaturated domain fails.

For group coefficients, the invisible coherent groups \(H_U\) are transported by a multiplicative equivalence, and the parameterizations \(x_0h\) commute with pullback. Thus feasibility and ambiguity are both retained. This is more than a comparison of the number of extensions. It is not an assertion that every possible based gauge orbit set is presentation independent.

\subsection{Why a boundary must retain its identification}
Let \(R:\mathcal C\to\mathcal D\) be a functor between groupoids, and fix \(b\in\mathcal D\). A realization over \(b\) should include a chosen isomorphism to \(b\), rather than require literal equality of representatives.

\begin{definition}[Marked realization groupoid]
The objects of \(\mathcal E_b(R)\) are pairs \((x,\phi)\), with \(x\in\mathcal C\) and \(\phi:Rx\xrightarrow{\sim}b\). An arrow \((x,\phi)\to(y,\psi)\) is an arrow \(a:x\to y\) with \(\psi R(a)=\phi\).
\end{definition}
Composition and inverses preserve this equation, so this is a groupoid. It is the explicit groupoid homotopy fiber construction, without invoking a higher categorical interface.

\begin{proposition}[Existence and invisible automorphisms]\label{KS38}
The groupoid \(\mathcal E_b(R)\) is nonempty exactly when \(b\) is in the essential image of \(R\). An automorphism of \((x,\phi)\) is an automorphism \(a\) of \(x\) satisfying \(R(a)=1_{Rx}\).
\end{proposition}
\begin{proof}
The first assertion is the object definition. For the second, the boundary equation is \(\phi R(a)=\phi\). Cancel the isomorphism \(\phi\). Conversely \(R(a)=1\) makes the equation hold.
\end{proof}
\source{KS38}
This retains precisely the symmetries invisible at the identified boundary. It avoids freezing arbitrary representative coordinates and then mistaking the result for an invariant realization problem.

\plate{08_marked}{The chosen boundary identification is part of a realization. A comparison must commute with it. Under an equivalent presentation, the comparison square may commute through an isomorphism rather than a literal equality, and the marked realization groupoid is still transported correctly.}{fig:marked}

\subsection{Weakly commuting squares are enough}
Suppose \(E:\mathcal C\to\mathcal C'\) and \(K:\mathcal D\to\mathcal D'\) are equivalences, with functors \(R,R'\) and a natural isomorphism
\[
 \sigma:KR\xRightarrow{\sim}R'E.
\]
The square is not required to commute literally.

\begin{theorem}[Invariance of the marked realization problem]\label{KS39}
The map
\[
 (x,\phi)\longmapsto(Ex,K(\phi)\sigma_x^{-1}),\qquad a\longmapsto E(a)
\]
defines an equivalence \(\mathcal E_b(R)\simeq\mathcal E_{Kb}(R')\).
\end{theorem}
\begin{proof}
Naturality of \(\sigma\) verifies the boundary equation for transported arrows. Full faithfulness follows by lifting an arrow through the fully faithful \(E\), then reflecting its boundary equation through the faithful \(K\).

For essential surjectivity, start from \((y,\psi)\). Choose \(x\) and an isomorphism \(i:Ex\to y\). The composite
\[
 \psi\,R'(i)\,\sigma_x:KRx\to Kb
\]
lifts through full \(K\) to an arrow \(\phi:Rx\to b\). Since the target is a groupoid it is an isomorphism. Its transported identification is \(\psi R'(i)\), so \(i\) gives the required marked comparison. The functor is fully faithful and essentially surjective, hence an equivalence.
\end{proof}
\source{KS39}
Apply this to the actual restriction functors between cocycle patch groupoids and the equivalences induced by an equivalent view presentation. The full realization problem, including comparison witnesses and automorphisms, is transported.

Replacing \(b\) by an isomorphic boundary \(b'\) also gives an equivalence, by composing the identification with \(b\to b'\). For \(R=\id\), the marked realization groupoid has center \((b,1_b)\), and between \((x,\phi)\) and \((y,\psi)\) there is exactly one compatible arrow \(\psi^{-1}\phi\). This is a useful example of contractible realization data, not a claim that the original groupoid has no symmetries.

